% E0 track: the adjoint on-ramp course document.
% Companion to tutorials/E_differentiable/{E0a,E0b,E0c}. Self-contained:
% only amsmath/amssymb/booktabs/hyperref. Compile: tectonic (x1 sufficient).
\documentclass[11pt]{article}
\usepackage{amsmath}
\usepackage{amssymb}
\usepackage{booktabs}
\usepackage{hyperref}

\setlength{\parskip}{0.4em}
\setlength{\parindent}{0pt}

% ---- lightweight remark boxes (no external packages) --------------------
\newcounter{remark}
\newenvironment{remark}[1]{%
  \refstepcounter{remark}\par\medskip
  \noindent\begin{minipage}{\linewidth}
  \hrule height 0.8pt \smallskip
  \textbf{Remark \theremark\ (#1).}\ }{%
  \smallskip\hrule height 0.8pt \end{minipage}\par\medskip}

\newenvironment{measured}{%
  \par\medskip\noindent\begin{minipage}{\linewidth}
  \hrule height 0.8pt \smallskip
  \textbf{Measured.}\ }{%
  \smallskip\hrule height 0.8pt \end{minipage}\par\medskip}

% ---- math shorthands ----------------------------------------------------
\newcommand{\diff}[2]{\frac{\partial #1}{\partial #2}}
\newcommand{\vd}[2]{\frac{\delta #1}{\delta #2}}
\newcommand{\ip}[2]{\left(#1,\,#2\right)}
\newcommand{\dR}{\,\mathrm{d}}
\newcommand{\RR}{\mathbb{R}}
\newcommand{\nn}{\mathbf{n}}

\title{The Differentiable Simulation Track (E0--E2):\\
  Mathematical Companion\\
  \large E0 = the adjoint on-ramp; E1/E2 = the payoff chapters}
\author{The DiffSim curriculum}
\date{July 2026}

\begin{document}
\maketitle

\begin{abstract}
This document is the theory companion to the full E-track of differentiable
simulation tutorials (\texttt{tutorials/E\_differentiable/E0a--E2}).
\textbf{E0} (Parts I--VII, the on-ramp) develops the mathematics of
differentiable simulation in DiffSim notation: three paths to a gradient,
the discrete-adjoint stance, taping rules, transient chains, and $J$-design
principles.  Structure and example lineage are adapted from the mathematical
background of \textbf{dolfin-adjoint/pyadjoint} by Patrick E.~Farrell; the
mother problem, three-way comparison device, and ``learning reaction rates''
examples follow Farrell, Ham, Funke \& Rognes \cite{farrell2013} and
Mitusch, Funke \& Dokken \cite{mitusch2019}.  DiffSim-native content ---
the discrete-adjoint stance, taping rules, the beyond-FH gauge story, the
H4 signal-scale protocol --- is developed in the E-track specs and dev
notes.

\textbf{E1 and E2} (Part VIII, the payoff chapters) apply the on-ramp
machinery to geometric inverse problems: E1 recovers a hidden disk from
sparse Poisson probes by differentiating through the Shifted-Boundary
Method (SBM) face residuals and the implicit-function-theorem
closest-point projection; E2 connects a GENIE implicit neural
representation (INR) as the geometry oracle and recovers a hidden
last-layer edit from probe data alone, closing the differentiable
simulation $\times$ neural geometry loop.

House rule: \emph{theory predicts a number, then we print the measured
number next to it.}  All E0 numbers quoted here are from the gate batteries
of tutorials E0a--E0c, reproducible on CPU in under three minutes each.
E1/E2 numbers are from the respective chapters at the grid levels noted.
\end{abstract}

\tableofcontents

%=======================================================================
\section*{Part I: The Reduced Functional and Three Paths to Its Gradient}
\addcontentsline{toc}{section}{Part I: The Reduced Functional}

\section{The reduced functional}
\label{sec:reduced}

Let $m \in \RR^N$ be a vector of simulation parameters (a conductivity
field per element, material coefficients, boundary data, \ldots).  The
forward model maps $m$ to a state $u(m) \in \RR^n$ by solving a residual
equation
\begin{equation}
  R\bigl(u(m),\, m\bigr) \;=\; 0 .
  \label{eq:residual}
\end{equation}
A scalar quantity of interest $J(u, m)$ defines the \emph{reduced
functional}
\begin{equation}
  \hat{J}(m) \;=\; J\!\bigl(u(m),\, m\bigr),
  \label{eq:Jhat}
\end{equation}
where $u(m)$ is understood as the implicit solution of \eqref{eq:residual}.
The optimisation problem is $\min_m \hat{J}(m)$; gradient-based methods
require $\nabla_m \hat{J}$.

\textbf{The mother problem (Farrell/pyadjoint).}  Throughout Part~I we
use the Poisson conductivity problem as a concrete instance:
$-\nabla\cdot(\kappa(x)\nabla u) = f$ on $\Omega = [0,1]^2$ with
$u|_{\partial\Omega} = g$.  The conductivity field $\kappa(x)$ is our
parameter, discretized as a per-element constant $m = (\kappa_1,\ldots,
\kappa_N)$ on a uniform $Q_1$ mesh.  Someone hands us probe measurements
$u_{\mathrm{obs}}(x_i)$ at $P$ interior locations; the objective is
\begin{equation}
  J = \tfrac12 \sum_{i=1}^P \bigl(u(x_i) - u_{\mathrm{obs},i}\bigr)^2 .
  \label{eq:Jprobe}
\end{equation}
For $N = 64$ elements (level-3 grid) the tutorial measures
$\hat{J}(\kappa_{\mathrm{eval}}) \approx 8.35 \times 10^{-3}$ at a
uniform evaluation point $\kappa_{\mathrm{eval}} = 1.3$, with a two-blob
true field $\kappa_{\mathrm{true}} \in [1.0,\, 2.5]$.

\section{Three ways to $\nabla_m \hat{J}$}
\label{sec:threeways}

Differentiating the residual constraint \eqref{eq:residual}
implicitly links $\nabla_m \hat{J}$ to $u$ and $R$.  Three methods
exploit this link at radically different cost.

\subsection{Finite differences (FD)}

For each parameter index $i = 1,\ldots,N$:
\begin{equation}
  \frac{\partial \hat{J}}{\partial m_i}
  \;\approx\;
  \frac{\hat{J}(m + \epsilon e_i) - \hat{J}(m - \epsilon e_i)}{2\epsilon}.
  \label{eq:fd}
\end{equation}
Cost: \textbf{$2N$ forward solves} (central differences).  Accuracy:
limited by floating-point cancellation for $\epsilon \lesssim
\sqrt{\varepsilon_{\mathrm{machine}}}$.  At $N = 64$ this requires 128
forward solves.

\subsection{Tangent-linear method (TLM, forward-mode AD)}

Differentiate \eqref{eq:residual} with respect to $m_i$:
\begin{equation}
  \frac{\partial R}{\partial u} \,\dot{u}_i + \frac{\partial R}{\partial m_i}
  = 0
  \quad\Rightarrow\quad
  \dot{u}_i = -\!\left(\frac{\partial R}{\partial u}\right)^{\!-1}
              \frac{\partial R}{\partial m_i} .
  \label{eq:tlm}
\end{equation}
Each $\dot{u}_i$ is obtained by one \emph{forward-mode} linear solve, then
$\partial \hat{J}/\partial m_i = (\partial J/\partial u)\,\dot{u}_i$.
Cost: \textbf{$N$ forward-mode solves} --- half the FD cost, same
$O(N)$ scaling.

\subsection{Adjoint (reverse-mode)}
\label{sec:adjoint}

Form the Lagrangian by appending the residual constraint with a multiplier
$\lambda$:
\begin{equation}
  \mathcal{L}(u, m, \lambda)
  \;=\; J(u, m) \;-\; \lambda^\top R(u, m) .
  \label{eq:lagrangian}
\end{equation}
Stationarity with respect to $u$ gives the \textbf{adjoint equation}:
\begin{equation}
  \frac{\partial \mathcal{L}}{\partial u} = 0
  \quad\Longrightarrow\quad
  \left(\frac{\partial R}{\partial u}\right)^{\!\top} \lambda
  = \left(\frac{\partial J}{\partial u}\right)^{\!\top} .
  \label{eq:adj}
\end{equation}
This is \emph{one linear system for $\lambda$, independent of $N$}.
Once $\lambda$ is known, the gradient with respect to every parameter
follows from $\partial \mathcal{L}/\partial m = 0$:
\begin{equation}
  \frac{\partial \hat{J}}{\partial m}
  = \frac{\partial J}{\partial m}
  - \lambda^\top \frac{\partial R}{\partial m} .
  \label{eq:adjgrad}
\end{equation}
Cost: \textbf{1 forward solve $+$ 1 adjoint solve}, regardless of $N$.

\begin{remark}{adjoint linearity}
Even when $R(u,m) = 0$ is nonlinear in $u$, the adjoint equation
\eqref{eq:adj} is \emph{always linear} in $\lambda$ --- it is the
transposed Jacobian of $R$ evaluated at the converged state $u^*(m)$.
There is no iteration, only one linear solve.
\end{remark}

\begin{table}[h]
\centering
\caption{Three paths to $\nabla_m \hat{J}$, cost in forward-solve units.
At $N = 64$ the tutorial measures an FD/adjoint wall-time ratio of
approximately 30--80$\times$ (kernel compilation dominates the first run;
steady-state ratio $\approx 35\times$).}
\label{tab:threeways}
\begin{tabular}{lccc}
\toprule
Method & Solve cost & Exact? & Scales with $N$? \\
\midrule
Finite differences (central) & $2N$ forward & No ($\epsilon$-limited) & Yes \\
Tangent-linear                & $N$ forward  & Yes                    & Yes \\
Adjoint                       & $1 + 1$      & Yes                    & \textbf{No} \\
\bottomrule
\end{tabular}
\end{table}

For the Poisson mother problem with $N = 64$ the adjoint equation is
$A^\top \lambda = (\partial J/\partial u)^\top$, where $A(\kappa)$ is
the assembled constrained stiffness.  Since $A$ is symmetric
($A^\top = A$), the forward-solve LU factorization can be reused
directly for the adjoint solve at no additional cost.

\begin{measured}
Level-3 grid, 64 elements, 5 probes, FP64.  Adjoint vs central FD:
max relative error $< 10^{-6}$.  Adjoint vs \texttt{wp.Tape} (taped
in $\kappa$): max relative error $< 10^{-9}$.  FD wall time / adjoint
wall time: 30--80$\times$ (first run); $\approx 35\times$ steady-state.
(\texttt{tutorials/E\_differentiable/E0a\_thinking\_differentiable.py}.)
\end{measured}

\section{Information-flow reversal}
\label{sec:infoflow}

The forward solve propagates information \emph{from} the parameter field
$m$ \emph{through} the PDE \emph{to} the output functional $J$.  The
adjoint reverses this flow: it carries sensitivity \emph{from} the scalar
$J$ \emph{backward through} the PDE \emph{to} every parameter at once.
This is why the adjoint cost is $O(1)$ in $N$: instead of differentiating
$N$ forward paths, it differentiates one backward path.

The information-flow picture is due to Farrell~\cite{farrell2013} and is
the central image of the E0a tutorial.

%=======================================================================
\section*{Part II: Discretize-Then-Differentiate}
\addcontentsline{toc}{section}{Part II: Discretize-Then-Differentiate}

\section{The discrete-adjoint stance}
\label{sec:dtd}

Two philosophies exist for adjoints of PDE-constrained problems.

\emph{Differentiate-then-discretize} derives the adjoint PDE in the
continuous setting (a second boundary-value problem) and then discretizes
it.  The resulting discrete adjoint is not the exact transpose of the
primal discretization; it introduces an additional error term of the same
order as the spatial discretization.

\emph{Discretize-then-differentiate} first commits to a discrete residual
$R_h(u_h, m) = 0$ and then differentiates the discrete system.  The
resulting discrete adjoint satisfies
$({\partial R_h}/{\partial u_h})^\top \lambda = ({\partial J_h}/
{\partial u_h})^\top$ \emph{exactly}---there is no extra approximation
beyond the chosen discretization.

\textbf{DiffSim's stance is discretize-then-differentiate} (spec S4.3).
Every formula in these notes applies to the discrete system; the word
``adjoint equation'' always refers to the discrete transpose.

\begin{remark}{IFT / relation vs.\ algorithm}
The Implicit Function Theorem (IFT) underlies the entire approach.
Given a smooth residual $R(u, m) = 0$ with $\partial R/\partial u$
invertible at a solution $(u^*, m)$, the IFT guarantees a locally smooth
map $m \mapsto u^*(m)$ whose derivative is
\begin{equation}
  \frac{\partial u^*}{\partial m}
  = -\!\left(\frac{\partial R}{\partial u}\right)^{\!-1}
    \frac{\partial R}{\partial m}\biggr|_{(u^*,m)} .
  \label{eq:ift}
\end{equation}
We differentiate the \emph{relation} $R = 0$, not the \emph{algorithm}
used to solve it.  A Newton iteration, a Krylov method, or an LU
factorization is an algorithm for inverting $R$; none of these need to
be differentiated.  Only the assembled matrix $\partial R/\partial u$
evaluated at convergence enters the adjoint equation.
\end{remark}

\section{The Lagrangian derivation of the discrete adjoint}
\label{sec:lagrangian}

Let $R_h : \RR^n \times \RR^N \to \RR^n$ be the discrete residual and
$J_h : \RR^n \times \RR^N \to \RR$ the discrete objective.
The Lagrangian is
\begin{equation}
  \mathcal{L}(u_h, m, \lambda) = J_h(u_h, m) - \lambda^\top R_h(u_h, m) .
\end{equation}
Stationarity conditions:
\begin{align}
  \frac{\partial \mathcal{L}}{\partial \lambda} &= 0
  \quad\Longrightarrow\quad R_h(u_h^*, m) = 0
  \quad\text{(forward problem)}, \label{eq:stat_lam}\\
  \frac{\partial \mathcal{L}}{\partial u_h} &= 0
  \quad\Longrightarrow\quad
  \underbrace{\left(\frac{\partial R_h}{\partial u_h}\right)^{\!\top}}_{A^\top}
  \lambda
  = \underbrace{\left(\frac{\partial J_h}{\partial u_h}\right)^{\!\top}}_{\text{RHS}}
  \quad\text{(adjoint equation)}, \label{eq:stat_u}\\
  \frac{\partial \mathcal{L}}{\partial m} &= 0
  \quad\Longrightarrow\quad
  \frac{\partial \hat{J}}{\partial m}
  = \frac{\partial J_h}{\partial m}
  - \lambda^\top \frac{\partial R_h}{\partial m}
  \quad\text{(gradient formula)}. \label{eq:stat_m}
\end{align}
For the Poisson problem, $\partial R_h/\partial u_h = A(\kappa)$ and
$\partial R_h/\partial \kappa_e = A_e$ (the element stiffness at
$\kappa = 1$), so \eqref{eq:stat_m} reduces to
\begin{equation}
  \frac{\partial \hat{J}}{\partial \kappa_e}
  = -\lambda_{\mathrm{int}}^\top A_e\, u_h^*,
  \label{eq:poissongrad}
\end{equation}
where $\lambda_{\mathrm{int}}$ is $\lambda$ with boundary-node
entries zeroed (because the Dirichlet-eliminated rows of $A$ do not
depend on $\kappa_e$).

\begin{remark}{symmetry and factorization reuse}
For the Poisson operator $A = A^\top$, so the adjoint equation
$A^\top \lambda = \partial J_h/\partial u_h$ is solved with the
\emph{same LU factorization} as the forward solve.  In
\texttt{scipy.sparse.linalg.splu}, the call
\texttt{lu.solve(rhs,~trans=`T')} triggers the transposed
forward-backward substitution using the stored triangular factors.
For non-symmetric operators (SBM, convection-diffusion, Navier-Stokes)
the same principle holds: one factorization, two solve directions.
\end{remark}

%=======================================================================
\section*{Part III: Anatomy of a Taped Brick}
\addcontentsline{toc}{section}{Part III: Anatomy of a Taped Brick}

\section{What \texttt{wp.Tape} records}
\label{sec:tape}

A Warp \texttt{wp.Tape} is a \emph{launch graph}: a list of
\[
  (\text{kernel\_ref},\; \text{dim},\; \text{input arrays},\;
   \text{output arrays})
\]
tuples appended during the \texttt{with tape:} scope.  On
\texttt{tape.backward()}, the list is replayed in reverse; for each entry
the backward kernel propagates cotangents
\[
  \mathrm{adj}\_x \mathrel{+}= \frac{\partial f}{\partial x}\,
  \mathrm{adj}\_y,
\]
where $\mathrm{adj}\_y = \texttt{tape.gradients}[y]$ is the incoming seed.
Arrays with \texttt{requires\_grad=False} are \emph{frozen}: the tape
skips cotangent accumulation for them.

\textbf{Custom VJP vs.\ unrolled tape.}  When a kernel has
\texttt{enable\_backward=False} (e.g.\ the Poisson assembly kernel in
\texttt{diffsim/assembly/}\allowbreak\texttt{operators.py}), the tape
does \emph{not} generate an automatic backward kernel.  The gradient path
is instead supplied as a \emph{custom VJP} --- in DiffSim's case, the
transposed solve $A^\top \lambda = \partial J/\partial u$.  This is
intentional: letting the tape unroll an LU factorization would be
catastrophically expensive and unnecessary, given the IFT.

\section{The do-not-differentiate list}
\label{sec:dontdiff}

Each item below is an \emph{algorithm} implementing a mathematical
\emph{relation}.  The relation has a clean adjoint; the algorithm's
code graph does not.

\begin{enumerate}
\item \textbf{Linear solvers (splu, bicgstab, GMRES, cuDSS, \ldots).}
  Relation: $A(m)\,u = b$.  Adjoint: $A^\top \lambda = \partial J/
  \partial u$ (one solve, same LU factors).  Never unroll the
  LU/Krylov steps through a tape.

\item \textbf{Newton / nonlinear iterations.}
  Relation: $R(u^*, m) = 0$ defines $u^*(m)$ implicitly (IFT).
  Adjoint: $(\partial R/\partial u)^\top \lambda = \partial J/\partial u$
  evaluated at the \emph{converged} state.  Never tape through the
  iteration loop.

\item \textbf{Mesh build, octree carve, element classification,
  preflight checks.}
  These produce combinatorial objects (connectivity tables, sorted
  element lists) that are piecewise constant in the geometric parameters;
  their derivative is zero almost everywhere and undefined at transitions.
  The correct approach: \emph{freeze} the mesh before opening the tape.
  Geometric sensitivities live in the boundary parametrization (shape
  derivatives), not the mesh arrays.  E1
  (\texttt{E1\_shape\_optimization.py}) is the worked instance.

\item \textbf{RNG seeds and tabulated basis functions.}
  Basis tables (\texttt{N}, \texttt{dN}, \texttt{w}) are computed once by
  \texttt{basis\_tables(p,~dim)} and treated as frozen inputs with
  \texttt{requires\_grad=False}.  Random seeds are inputs in the
  reparametrization sense; they are never inside the tape scope.
\end{enumerate}

\section{Compute-once-and-store patterns}
\label{sec:computeonce}

Three canonical patterns avoid redundant computation.

\textbf{(i) One LU factorization for $A$ and $A^\top$ solves.}
\texttt{scipy\-.sparse\-.linalg\-.splu} stores both triangular factors; the
forward solve (\texttt{lu.solve(b)}) and the adjoint solve
(\texttt{lu.solve(c,~trans=`T')}) share the same factorization.

\begin{measured}
Level-3 Poisson, 50 trials.  One LU (two solve calls): $\approx t$~ms.
Two separate LU factorizations: $\approx 1.7t$~ms.  Ratio $> 1.2\times$
(measured; the factorization cost dominates at this size; the ratio
grows toward $2\times$ at scale because factorization is
$O(n^{1.5})$ in 2-D while each solve is $O(n\log n)$).
(\texttt{E0b}, \S4.)
\end{measured}

\textbf{(ii) Tabulated basis reuse.}  \texttt{dm.bins[p][`dN']} is
computed once; every kernel launch takes a reference to the same device
array.  Its shape is \texttt{(nqp, nbf, dim)} --- at $Q_1$ on a 2-D
level-3 grid this is \texttt{(4, 4, 2)}.

\textbf{(iii) GP-field precomputation (\emph{frozen at GPs}).}  A
closure field $\kappa_\theta(x)$ is evaluated at all Gauss points
\emph{before} the tape scope, stored in an array with
\mbox{\texttt{requires\_grad=False}}, and passed as a frozen input
to the assembly kernel.  The tape tracks only the PDE residual's
dependence on the tabulated values, not the internals of
$\kappa_\theta$.  This is the M2 ``one-way coupler'' convention:
the closure changes in the outer optimisation loop; the PDE adjoint
needs only $\partial R/\partial \kappa$ at the quadrature points.

\section{The dot-product test}
\label{sec:dotproduct}

For any operator $A$ and vectors $v$, $w$:
\begin{equation}
  \langle A v,\, w \rangle \;=\; \langle v,\, A^\top w \rangle .
  \label{eq:dottest}
\end{equation}
This is the cheapest possible sanity check for a new operator and its
adjoint.  The test does not require an explicit matrix: only the
forward action $v \mapsto Av$ and the adjoint action $w \mapsto A^\top w$
are needed (it is valid for matrix-free operators).  A failure means the
adjoint is wrong; a pass to machine precision gives strong evidence it
is correct.

The verification protocol for every new DiffSim operator: pick random
$v, w$; compute $Av$ and $A^\top w$; check
\[
  \frac{|\langle Av,w\rangle - \langle v,A^\top w\rangle|}
  {\|Av\|\,\|w\| + 10^{-100}}
  \;<\; 10^{-12} .
\]

\begin{measured}
Level-3 Poisson assembled CSR.  Relative error of dot-product test:
$\sim 10^{-17}$ (Poisson is symmetric, $A = A^\top$, so the test is
algebraically exact up to floating-point rounding).
(\texttt{E0b}, \S5.)
\end{measured}

%=======================================================================
\section*{Part IV: Transient Chains and Compute-vs-Memory}
\addcontentsline{toc}{section}{Part IV: Transient Chains}

\section{The transient adjoint}
\label{sec:transient}

A time-marching scheme produces a sequence of states
$u^0, u^1, \ldots, u^N$.  Each step is a residual
\begin{equation}
  F^n(u^n, u^{n-1}, m) = 0 , \qquad n = 1, \ldots, N .
\end{equation}
For BDF1 heat: $F^n = A\,u^n - (M/\Delta t)\,u^{n-1} = 0$, with
$A = M/\Delta t + K(\kappa)$.  The Lagrangian appends one multiplier
$\lambda^n$ per step:
\begin{equation}
  \mathcal{L} = J_h(u^N, m)
  - \sum_{n=1}^N (\lambda^n)^\top F^n(u^n, u^{n-1}, m).
\end{equation}
Stationarity at $u^n$ for $n < N$:
\begin{align}
  \left(\frac{\partial F^n}{\partial u^n}\right)^\top \lambda^n
  &= -\left(\frac{\partial F^{n+1}}{\partial u^n}\right)^\top \lambda^{n+1}
  \nonumber\\
  A^\top \lambda^n &= P^\top \lambda^{n+1},
  \qquad P = M/\Delta t , \label{eq:adjrec}
\end{align}
with terminal condition at $n = N$:
\begin{equation}
  A^\top \lambda^N = \left(\frac{\partial J_h}{\partial u^N}\right)^\top
  = W^\top(W u^N - u_{\mathrm{obs}}) .
  \label{eq:adjterm}
\end{equation}
This is the \emph{adjoint recurrence}: solve backward in time, one
transposed solve per step.  The gradient contribution from step $n$ is
\begin{equation}
  \Delta\!\left(\frac{\partial \hat{J}}{\partial \kappa}\right)_e
  \mathrel{+}=
  -(\lambda^n_{\mathrm{int}})^\top A_e\, u^n,
  \label{eq:trangrad}
\end{equation}
where $u^n$ is the \emph{primal state stored from the forward march}
(checklist item~3 in Appendix~\ref{app:checklist}).

\begin{remark}{one factorization, all steps}
When $\kappa$ is constant in time, $A = M/\Delta t + K(\kappa)$ does
not change between steps.  A \emph{single LU factorization} serves every
forward solve and every adjoint solve --- forward via \texttt{lu.solve(rhs)}
and backward via \texttt{lu.solve(seed,~trans=`T')}.  Both directions
reuse the same stored triangular factors.
\end{remark}

\begin{measured}
BDF1 heat, level-3 grid, 30 steps, $\Delta t = 0.01$.  Transient-chain
check (discrete adjoint vs central FD, end-to-end): max relative error
$< 10^{-5}$.  (\texttt{E0c}, \S2.)
\end{measured}

\section{Compute-vs-memory: full-store and recompute}
\label{sec:cvm}

The backward sweep \eqref{eq:trangrad} requires $u^n$ at each step.  Two
strategies supply it.

\textbf{Full-store.} Keep $u^0, u^1, \ldots, u^N$ in memory.  Memory:
$(N+1) \times n_{\mathrm{nodes}} \times 8$~bytes (FP64).  Backward cost:
$N$ transposed solves.  At $10^6$ dofs and $10^3$ steps this is 8~GB
--- a regime where full-store exceeds GPU memory.

\textbf{Recompute from checkpoints.}  Store $n_c$ checkpoints at
intervals of $\Delta n = N/n_c$ steps.  When the backward sweep needs
$u^n$, re-march forward from the nearest earlier checkpoint.  Memory:
$n_c \times n_{\mathrm{nodes}} \times 8$~bytes.  Backward cost: $N$
transposed solves plus $O(n_c \cdot \Delta n)$ extra forward solves.

The optimal checkpoint schedule is \emph{revolve} (Griewank \& Walther
\cite{griewank2000}): $O(\log N)$ checkpoints achieve $O(N \log N)$
recompute at $O(\log N)$ memory, provably optimal.  Binomial scheduling
libraries (\texttt{pyrevolve}, \texttt{checkpoint\_schedules}) implement
this and wire directly into the backward loop above.

\textbf{The $2.5\times$ cost contract.}  A correct reverse-mode adjoint
adds at most one transposed solve plus one gradient contraction per forward
step.  The backward march should cost $\leq 2.5\times$ the forward march.
If backward $\gg$ forward, something is being unrolled that should not be.

\begin{measured}
BDF1 heat, level-3, 30 steps.  Full-store backward / forward wall-time
ratio $\leq 2.5\times$ (measured).  Primal-storage inventory: 31 states
$\times$ 81 nodes $\times$ 8 bytes = 20~KiB.  Full-store vs recompute-
from-6-checkpoints gradient agreement: $< 10^{-10}$ (same discrete
adjoint, bit-for-bit identical $u^n$ from re-march).
(\texttt{E0c}, \S3.)
\end{measured}

%=======================================================================
\section*{Part V: Choosing a Good $J$}
\addcontentsline{toc}{section}{Part V: Choosing a Good $J$}

\section{Smoothness: nonsmooth $J$ and its relaxation}
\label{sec:smoothness}

Functionals built on $\max$, $\min$, $|\cdot|$, or hard thresholds are
nonsmooth: the gradient jumps when the active branch switches, and is
zero almost everywhere for a threshold.  An optimiser fed a jumpy
gradient stalls or oscillates.

The standard cure is a \emph{relaxed surrogate}.  For
$J = \max_i g_i(p)$, the logsumexp (softmax) upper bound
\begin{equation}
  J_\beta(p) = \frac{1}{\beta}\log\sum_i \exp\!\bigl(\beta\, g_i(p)\bigr)
  \;\xrightarrow{\beta\to\infty}\; \max_i g_i(p)
  \label{eq:lse}
\end{equation}
has a smooth softmax-weighted gradient at any finite $\beta$.

\begin{measured}
Sweeping a scalar $p$ over $[0,3]$: max-J discrete gradient maximum jump
$\approx 0.55$; logsumexp ($\beta = 40$) maximum jump $\approx 0.17$;
ratio $\approx 0.31 < 0.5$ (smoother by $1.8\times$).
(\texttt{E0c}, \S4(i).)
\end{measured}

\section{Signal scale: the H4 protocol}
\label{sec:signalscale}

Before launching any optimiser, verify that $|\partial \hat{J}/\partial m|$
clears the finite-difference noise floor.  The \emph{H4 protocol}
(DiffSim-native): compute central FD across a range of $\epsilon$ values
and take the minimum FD spread as the noise-floor estimate.  If the adjoint
signal does not exceed the floor by orders of magnitude, the objective
barely responds to the parameters --- rescale $m$ or $J$, or choose a
more informative observable, before spending a single optimiser iteration.

\begin{measured}
Heat chain, level-3.  Adjoint signal $|\partial \hat{J}/\partial \kappa|_\infty
\approx 10^{-3}$; FD noise floor (min spread across $\epsilon \in
[10^{-3}, 10^{-8}]$) $\approx 10^{-7}$; signal-to-floor ratio
$\approx 10^4 \gg 10^3$ (required).
(\texttt{E0c}, \S4(ii).)
\end{measured}

\section{Flat-$J$ geometry and informative probe placement}
\label{sec:flatJ}

A probe placed where the solution field has negligible amplitude sees
$\partial u/\partial m \approx 0$: the gradient is near zero regardless of
how good the adjoint is.  The observable placement is part of the
experiment design, not the solver.

\begin{measured}
Heat chain, sharp bump at $(0.25, 0.25)$.  Probe on the active bump:
$\|\partial \hat{J}/\partial\kappa\|_\infty \approx 2 \times 10^{-3}$.
Probe in the far corner $(0.84, 0.84)$: $\approx 2 \times 10^{-5}$.
Contrast $\approx 100\times \gg 10\times$ (required).
(\texttt{E0c}, \S4(iii).)
\end{measured}

\section{Identifiability, gauge freedom, and regularisation as science}
\label{sec:identifiability}

A gradient can be correct, large, and smooth while the inverse problem
remains ill-posed: the parameterisation may have \emph{unidentifiable
directions} --- parameter combinations that produce zero observable change.
The diagnostic is the spectrum of the observation Jacobian (the Gramian
$G = J_m^\top J_m$ where $J_m = \partial \mathbf{d}/\partial m$ maps
parameters to data), not the gradient of $J$ itself.

\textbf{Case study: the beyond-FH gauge story (M4, rung~2).}
Learning a free-energy correction $f'_\theta(c)$ on top of Flory-Huggins
has two structurally invisible modes.  A constant in $f'$ shifts the
chemical potential $\mu$ by a constant, which is invisible to conserved-$c$
dynamics.  The linear mode proportional to $(1 - 2c)$ is \emph{exactly
aliased} with the FH chi term: the trajectory difference along this
direction was measured at $5 \times 10^{-15}$.  These are \emph{gauge
freedoms}; they must be removed from the parameterisation before fitting.

Quantitatively (numbers quoted from the beyond-FH dev note
\cite{beyondfh2026}):
\begin{itemize}
\item Un-anchored basis $\{\chi, P_1, P_2, P_3\}$:
  observation-Jacobian condition number $\approx 1.4 \times 10^{16}$
  ($\chi$ and $P_1$ aliased --- effectively singular).
\item Anchored basis $\{\chi, P_2, P_3\}$ (Legendre modes orthogonal
  to $\{1, c\}$, removing the gauge by construction):
  condition number $\approx 18$ (well-posed).  Recovery of coefficients
  from a single interior trajectory: relative error $\approx 5.6 \times
  10^{-8}$.
\item Conditioning improvement: $\approx 7.6 \times 10^{14}$.
\end{itemize}
Even with a gauge-anchored parameterisation, a \emph{single narrow
trajectory} (visited composition range $0.46$--$0.56$) leaves higher
modes weakly constrained (condition number $\approx 2 \times 10^2$--$5
\times 10^2$).  A \emph{composition-diverse} protocol (range
$0.21$--$0.85$) reduces the condition number to $\approx 14$--$18$
($15$--$30\times$ improvement) and recovers $\{P_2, \ldots, P_5\}$ to
coefficient error $\approx 4 \times 10^{-5}$, while the single narrow fit
fails (coefficient error $\approx 0.9$--$1.0$, a $> 2 \times 10^4\times$
recovery gap).

\textbf{Tikhonov regularisation as a prior, not a hack.}  A penalty
$\varepsilon \|\theta\|^2$ encodes the hypothesis that the correction is
small.  In the beyond-FH fit it does not merely stabilise: it pulls the
ill-conditioned single-trajectory fit to a definite (wrong) answer rather
than to NaN, making the failure legible.  Regularisation states a
scientific hypothesis; the condition number of the observation Jacobian
tells you whether the data can test it.

\begin{remark}{the rule taught three times}
E2 (the recovery plateau), M2--M3 (the observability ladder), and the
beyond-FH study all arrive at the same principle: \emph{mechanism before
optimiser}.  When a fit stalls, measure identifiability before tuning
the descent.
\end{remark}

\section{Multi-observable $J$ and the SP-1 bridge}
\label{sec:multiobs}

A single scalar probe misfit is often flat or aliased (\S\S\ref{sec:flatJ}
and \ref{sec:identifiability}).  The cure is a richer,
\emph{instrument-space} observable --- or several:
\begin{equation}
  J = \sum_k w_k \| \mathcal{O}_k(u) - d_k \|^2 ,
  \label{eq:multiJ}
\end{equation}
where $\mathcal{O}_k$ are diverse operators (probe values, a structure
factor $S(q,t)$, a film height $h(t)$, a PSF-convolved image).  Diverse
protocols illuminate different directions of parameter space, turning the
underdetermined inverse of E0b (5 probes for 64 unknowns) into a
well-posed one.  The beyond-FH instrument-space $S(q)$ misfit backpropagates
to the closure coefficients matching FD to $1.8 \times 10^{-11}$
(structure factor computed as a differentiable Warp kernel).  SP-1 (R3)
builds the multi-observable, multi-protocol inversion on the production
stack.

%=======================================================================
\section*{Part VI: The Mini Phase-Field Payoff}
\addcontentsline{toc}{section}{Part VI: Mini Phase-Field Payoff}

\section{Allen-Cahn: learning mobility and gradient-energy coefficient}
\label{sec:acpayoff}

The Allen-Cahn equation
\begin{equation}
  \partial_t c = -M \bigl[ f'(c) - \kappa\,\Delta c \bigr],
  \qquad f(c) = \tfrac14(c^2-1)^2 ,
  \label{eq:ac}
\end{equation}
is the group's own physics (F-track, M4) and serves as the payoff example
for the E0c recipe.  The ``learning reaction rates'' framing is the
dolfin-adjoint tutorial of that name \cite{mitusch2019}; we adapt it to
the Warp discrete-adjoint stack.

We differentiate $J = \tfrac12 \|W c^N - c_{\mathrm{obs}}\|^2$ with
respect to the scalar mobility $M$ and the gradient-energy coefficient
$\kappa$ through a 25-step semi-implicit BDF1 march.  The adjoint
recurrence \eqref{eq:adjrec} takes the form
\begin{align}
  A^\top \lambda^n &= \mathrm{seed}^n, \label{eq:acadjrec}\\
  \mathrm{seed}^{n-1} &= -\frac{\partial F^n}{\partial c^{n-1}}\,\lambda^n
  = \left[\frac{m_\ell}{\Delta t} - M\, m_\ell\,(3(c^{n-1})^2 - 1)\right]\lambda^n,
  \label{eq:acseed}
\end{align}
where $m_\ell$ is the lumped mass diagonal.

The three-way gradient check uses central FD as one leg and
\emph{complex-step differentiation} as the second exact leg: for a
holomorphic extension $\hat{J}(m + i\,h)$ with $h = 10^{-30}$,
\begin{equation}
  \frac{\partial \hat{J}}{\partial m} = \frac{\operatorname{Im}
  \hat{J}(m + ih)}{h} + O(h^2) ,
  \label{eq:cstep}
\end{equation}
with no subtraction error.  Agreement of all three legs to
$< 10^{-6}$ relative error provides strong evidence of correctness
independent of any single computational path.

\begin{measured}
Allen-Cahn, level-3, 25 steps, $\Delta t = 0.005$.  Three-way check:
adjoint vs complex-step vs FD, max relative error (worst of $\partial
\hat{J}/\partial M$ and $\partial \hat{J}/\partial \kappa$) $< 10^{-6}$.
(\texttt{E0c}, \S5.)
\end{measured}

%=======================================================================
\section*{Part VII: Numerical Results Summary}
\addcontentsline{toc}{section}{Part VII: Numerical Results Summary}

\section{Measured gates}
\label{sec:gates}

\begin{table}[h]
\centering
\caption{Measured gate results from tutorials E0a--E0c (CPU, FP64,
level-3 grid unless otherwise noted).  All reproducible with
\texttt{python tutorials/E\_differentiable/E0\{a,b,c\}*.py}.}
\label{tab:gates}
\begin{tabular}{llll}
\toprule
Chapter & Gate & Required & Measured \\
\midrule
E0a & adj vs FD max rel err           & $< 10^{-6}$  & $< 10^{-6}$ \\
E0a & tape vs FD max rel err          & $< 10^{-6}$  & $< 10^{-6}$ \\
E0a & adj vs tape max rel err         & $< 10^{-9}$  & $< 10^{-9}$ \\
E0a & FD / adjoint cost ratio         & 30--80$\times$ & $\sim$35$\times$ \\
\midrule
E0b & dot-product test rel err        & $< 10^{-12}$ & $\sim 10^{-17}$ \\
E0b & factorization reuse ratio       & $> 1.2\times$  & $> 1.2\times$ \\
E0b & misfit drop (40 GD steps)       & $> 100\times$  & $> 100\times$ \\
\midrule
E0c & heat-chain adj vs FD            & $< 10^{-5}$  & $< 10^{-5}$ \\
E0c & store vs recompute dJ agreement & $< 10^{-10}$ & $\sim 0$ \\
E0c & backward/forward wall ratio     & $< 2.5\times$  & $< 2.5\times$ \\
E0c & logsumexp/max gradient-jump ratio & $< 0.5$    & $\approx 0.31$ \\
E0c & signal / FD-floor ratio         & $> 10^3$     & $\approx 10^4$ \\
E0c & flat-J probe contrast           & $> 10\times$   & $\approx 100\times$ \\
E0c & AC three-way (adj/cstep/FD)     & $< 10^{-6}$  & $< 10^{-6}$ \\
\bottomrule
\end{tabular}
\end{table}

%=======================================================================
\section*{Part VIII: The Payoff Chapters — E1 and E2}
\addcontentsline{toc}{section}{Part VIII: The Payoff Chapters (E1 and E2)}

%-----------------------------------------------------------------------
\section{E1: Shape Optimization --- Finding the Hidden Circle}
\label{sec:e1}

E1 (\texttt{tutorials/E\_differentiable/\allowbreak E1\_shape\_optimization.py})
is the smallest closed instance of the loop that motivates all of DiffSim:
a quantity of interest $J$ defined on simulation output, differentiated with
respect to \emph{geometry} (not material fields), driving gradient descent
with the mesh re-carved from scratch every iteration.  The on-ramp (E0)
taught the adjoint for material parameters; E1 shows it working in the wild
for shape.

\subsection{The inverse problem}
\label{sec:e1problem}

A Poisson problem with conductivity $\kappa = 1.3$ is solved on a 2-D
square $[0,1]^2$ whose interior boundary is a disk of center
$(\theta_x^*, \theta_y^*)$ and radius $r^*$ (unknown).  Nine probe
values $u(x_i)$ measured on a small ring inside the disk are handed to the
recovery algorithm; the task is to find $\theta = (\theta_x, \theta_y, r)$
from these observations.

The objective is the standard probe least-squares misfit:
\begin{equation}
  J(\theta) \;=\; \tfrac{1}{2}
  \sum_{i=1}^{9} \bigl(u_\theta(x_i) - u^*_{\mathrm{obs},i}\bigr)^2,
  \label{eq:e1J}
\end{equation}
where $u_\theta$ is the FEM solution on the SBM mesh carved to the disk
parametrized by $\theta$.

\subsection{The gradient chain}
\label{sec:e1chain}

The full chain, as stated in the E1 docstring, is:
\begin{equation}
  J
  \;\xrightarrow{\partial J/\partial u}\;
  \lambda \;=\; (A^\top)^{-1}\!\left(\frac{\partial J}{\partial u}\right)^\top
  \;\xrightarrow{-\lambda^\top \partial R/\partial \theta}\;
  \theta .
  \label{eq:e1chain}
\end{equation}
Each arrow is one of the three payoff mechanisms; we expand them in order.

\subsubsection{Step 1: probe Jacobian $\partial J/\partial u$}

From \eqref{eq:e1J}, the right-hand side of the adjoint equation is
\begin{equation}
  \left(\frac{\partial J}{\partial u}\right)^\top
  \;=\; \sum_{i=1}^{9} \bigl(u_\theta(x_i) - u^*_{\mathrm{obs},i}\bigr)\,
  w_i,
  \label{eq:e1rhs}
\end{equation}
where $w_i$ are the point-evaluation weight vectors (sparse finite-element
interpolation weights) returned by \texttt{point\_eval\_weights}.
This is the same probe-misfit seed as in the E0 conductivity tutorial,
now assembled in the SBM geometry.

\subsubsection{Step 2: adjoint solve $A^\top \lambda = \partial J/\partial u$}

Nothing new relative to E0: one transposed LU solve.  Cost is one forward
solve, independent of the number of shape parameters (here 3, but the cost
is the same for 3000 GridSDF voxels --- the whole point of adjoints).  The
system matrix $A(\theta)$ is the SBM Nitsche-augmented stiffness for the
current geometry; it is assembled in the forward pass and its LU
factorization is reused for the adjoint (\S\ref{sec:computeonce}, item~i).

\subsubsection{Step 3: SBM face cotangents $-\lambda^\top \partial R/\partial \theta$}
\label{sec:e1sbmtape}

This is the step E0's do-not-differentiate list (\S\ref{sec:dontdiff},
item~3) pointed to as ``the worked instance.''  The SBM residual
$R(u, \theta) = 0$ depends on $\theta$ through:
\begin{itemize}
  \item the signed-distance vectors $\mathbf{d}(\theta)$ from each cut-face
    foot to the boundary (computed by the Newton closest-point projection),
  \item the mapped Dirichlet data $g(x + \mathbf{d}(\theta))$ evaluated at
    the projected boundary point.
\end{itemize}
A \texttt{wp.Tape} scope records the SBM face-residual kernels as
functions of these geometric quantities.  On \texttt{tape.backward()},
the tape propagates cotangents through the Nitsche face integrals back to
$\mathbf{d}$ and to $g(x + \mathbf{d})$:
\begin{equation}
  \frac{\partial R}{\partial \theta}
  \;=\; \frac{\partial R}{\partial \mathbf{d}}\,
  \frac{\partial \mathbf{d}}{\partial \theta}
  \;+\; \frac{\partial R}{\partial (g \circ p_\theta)}\,
  \frac{\partial (g \circ p_\theta)}{\partial \theta},
  \label{eq:e1dRdtheta}
\end{equation}
where $p_\theta(x)$ is the boundary projection map.  The tape delivers
the left-multiplied cotangent $-\lambda^\top (\partial R/\partial \mathbf{d})$
and $-\lambda^\top (\partial R/\partial g)$ as gradients on the intermediate
arrays.  These are exactly the ``tape cotangents for distance vectors +
mapped boundary data'' mentioned in the E1 docstring.

\subsubsection{Step 4: IFT through the Newton projection}
\label{sec:e1ift}

The distance vector $\mathbf{d}(\theta)$ is itself the output of a Newton
closest-point projection (a nonlinear iteration in \texttt{torch}) applied
to the oracle's signed-distance function $\psi(\cdot\,;\theta)$.
By the Implicit Function Theorem applied to the projection stationarity
condition,
\begin{equation}
  \frac{\partial \mathbf{d}}{\partial \theta}
  \;=\; -\!\left(\frac{\partial^2 \psi}{\partial x^2}\right)^{\!-1}
  \frac{\partial^2 \psi}{\partial x\, \partial \theta}\biggr|_{x = p(\theta)},
  \label{eq:e1ift}
\end{equation}
where $x = p(\theta)$ is the converged foot point.
\emph{We differentiate the projection relation, not the Newton
algorithm}---exactly the IFT stance of \S\ref{sec:dtd}, Remark~2.
In practice, \texttt{torch.autograd} differentiates through the oracle's
\texttt{psi(x; theta)} computation (a smooth CSG formula or a SIREN
network), and the Newton iteration is not unrolled.

This step closes the chain: the tape cotangent at $\mathbf{d}$ is
multiplied by $\partial \mathbf{d}/\partial \theta$ (via \texttt{torch}'s
backward) to deliver the final gradient $\partial J/\partial \theta$ on the
\texttt{torch} tensor \texttt{theta}.

\subsection{Why classification freezing makes the gradient well-defined}
\label{sec:e1freeze}

The E0 do-not-differentiate list (\S\ref{sec:dontdiff}, item~3) establishes
that mesh build, octree carve, and element classification are
\emph{piecewise constant} in geometric parameters: the classify step assigns
each cell a binary inside/outside label that does not change as $\theta$
moves smoothly within a classification epoch.  Therefore
\[
  \frac{\partial (\text{element classification})}{\partial \theta} = 0
  \quad\text{(almost everywhere in } \theta\text{).}
\]
E1 asserts this by tests.  The gradient is well-defined \emph{within} an
epoch because only the continuous geometric quantities ($\mathbf{d}$,
$g \circ p_\theta$) depend on $\theta$, not the combinatorial mesh
topology.  Between iterations the world is rebuilt: new retained set, new
surrogate, new constraints.  Gradients survive because they never depended
on mesh topology, only on the smooth geometric quantities.  This is the
``freeze-classify, differentiate-geometry'' principle that makes
shape-adjoint codes on cut meshes tractable.

\subsection{Connection to E0 concepts}
\label{sec:e1e0connect}

\begin{itemize}
  \item \textbf{IFT / relation vs.\ algorithm} (\S\ref{sec:dtd}): the Newton
    projection is an algorithm; only the converged foot-point relation
    enters \eqref{eq:e1ift}.
  \item \textbf{Do-not-differentiate list, item~3} (\S\ref{sec:dontdiff}):
    mesh build and classification are frozen before the tape opens.
  \item \textbf{Custom VJP for the linear solve} (\S\ref{sec:tape}): the
    adjoint solve $A^\top \lambda = \partial J/\partial u$ is a custom
    backward rule; the LU iteration is not unrolled.
  \item \textbf{Dot-product test} (\S\ref{sec:dotproduct}): the SBM
    Nitsche operator's adjoint is verified with the dot-product test before
    any shape optimization.
  \item \textbf{Signal scale / H4} (\S\ref{sec:signalscale}): the probe
    ring is placed inside the disk (not in the far field) to ensure
    $|\partial J/\partial \theta|$ clears the FD noise floor.
\end{itemize}

%-----------------------------------------------------------------------
\section{E2: GENIE/DiffSBM --- Neural Geometry Meets the Adjoint}
\label{sec:e2}

E2 (\texttt{tutorials/E\_differentiable/\allowbreak E2\_genie\_diffsbm.py})
extends E1's shape-gradient machinery to a provided implicit neural
representation (INR) as the geometry oracle.  The physical problem is
3-D SBM Poisson on a sphere-like surface; the design variables are the
INR's last-layer edit coefficients $\alpha \in \RR^k$.

\subsection{The provided-INR contract}
\label{sec:e2inr}

DiffSim does \emph{not} train INRs
(\texttt{docs/dev/specs/2026-07-06-neuralsdf.md}, \S N0).
A trained checkpoint is a \emph{provided input}: a SIREN-family network
with a linear last layer,
\begin{equation}
  \psi(x;\theta_L) \;=\; h_\phi(x)^\top \theta_L,
  \label{eq:inrpsi}
\end{equation}
where $h_\phi$ is the frozen feature extractor and $\theta_L$ is the
editable linear head.  DiffSim loads the checkpoint via
\texttt{ProvidedINROracle} and uses $\psi$ purely as an SDF oracle
(sign for classification, Newton+IFT for distance vectors).

\subsection{GENIE edit modes}
\label{sec:e2genie}

Following Karki et al.\ (GENIE) \cite{karki2025}, for linear-last-layer
INRs any last-layer edit $\Delta\theta_L$ changes the field by
$\Delta\psi(x) = h_\phi(x)^\top \Delta\theta_L$ --- the geometry is
\emph{affine} in the edit parameters.  The Gram operator over a
thick-band sampling distribution $\mu$ on a shell near the surface is
\begin{equation}
  G \;=\; \mathbb{E}_{x \sim \mu}\!\bigl[h_\phi(x)\, h_\phi(x)^\top\bigr],
  \label{eq:gram}
\end{equation}
and its top-$k$ eigenvectors $\{v_j\}_{j=1}^k$ define the \emph{design
space}: $\theta_L(\alpha) = \theta_L^0 + \sum_{j=1}^k \alpha_j v_j$.
Sampling on a \emph{thick} band (width $\approx \pm 0.05$ times the
surface radius, versus a thin shell of $\pm 0.005$) is essential for
stable modes: the GENIE paper's Fig.~4a finding is directly adopted as
DiffSim's sampling rule.  The mode stability is measured as the
subspace angle between two independent thick-band draws; at the sphere
checkpoint with $k=4$ modes, stability $= 1.000$ is measured.

\begin{remark}{DiffSim does not train INRs}
The program rule is: DiffSim consumes trained checkpoints but never
initiates a training run.  This boundary is an explicit design choice
(\texttt{docs/dev/specs/2026-07-06-neuralsdf.md}, \S N0): training INRs
is a different optimization problem (minimizing a neural rendering loss),
with different infrastructure, data pipelines, and convergence criteria.
The E2 chapter uses an analytical sphere checkpoint as the canonical test
double; real GENIE checkpoints drop in via \texttt{from\_genie\_json}.
\end{remark}

\subsection{Differentiating through the SDF query into the SBM machinery}
\label{sec:e2diffchain}

The INR gradient chain is structurally identical to E1's
(\S\ref{sec:e1chain}), with $\alpha$ replacing $\theta$:
\begin{equation}
  J
  \;\xrightarrow{\partial J/\partial u}\;
  \lambda \;=\; (A^\top)^{-1}(\partial J/\partial u)^\top
  \;\xrightarrow{-\lambda^\top \partial R/\partial \alpha}\;
  \alpha .
  \label{eq:e2chain}
\end{equation}

The new ingredient is the $\partial R/\partial \alpha$ path through the INR.
Because $\psi$ is affine in $\alpha$:
\begin{equation}
  \frac{\partial \psi}{\partial \alpha_j}(x)
  \;=\; h_\phi(x)^\top v_j,
  \label{eq:e2dpsi}
\end{equation}
which is smooth and closed-form (or inexpensive via \texttt{torch.autograd}
on the network forward).  The distance vector $\mathbf{d}(\alpha)$ is again
the output of the Newton closest-point projection applied to $\psi(\cdot;
\alpha)$, differentiated by the IFT (\S\ref{sec:e1ift} with $\psi$ now an
INR).  The tape cotangents through the SBM face residuals then deliver
$-\lambda^\top \partial R/\partial \alpha$ via \texttt{shape\_gradient},
and \texttt{torch.autograd} propagates through \eqref{eq:e2dpsi} to give
the final $\partial J/\partial \alpha$ on the \texttt{alpha} tensor.

\subsection{Epoch management: the trust-region contract}
\label{sec:e2epoch}

Because classification is frozen within an epoch (\S\ref{sec:e1freeze}),
the SBM stiffness $A(\alpha)$ is smooth in $\alpha$ only while the surface
stays within the frozen surrogate's trust region.  E2 defines the epoch
trust radius as
\begin{equation}
  \text{DRIFT} \;=\; 0.35\, h_{\min}, \qquad h_{\min} = 2^{-\text{LEVEL}},
  \label{eq:e2drift}
\end{equation}
where $h_{\min}$ is the finest-level cell size.  When
$\|\mathbf{d}(\alpha)\|_\infty > \text{DRIFT}$, a new classification epoch
is triggered at the current $\alpha$.  This \emph{epoch continuation}
strategy (also used in the M3 bunny hero, \texttt{docs/dev/roadmap.md})
gives smooth-within, valid-across: gradient descent can cross epoch
boundaries, with the objective continuous but the mesh topology changing
at each epoch transition.

\begin{remark}{branch stability and the FD check}
The three measured INR facts baked into the framework
(\texttt{docs/dev/specs/2026-07-06-neuralsdf.md}): the computational
WINDOW contract (feature spans $\geq 8$ cells at the chosen level), domain-
BOUNDED projections (sine INRs have spurious far zero sheets, so Newton is
initialized within the windowed sub-box), and the BRANCH-STABILITY rule:
at a bistable closest-point the projection selects different branches for
$\alpha + \epsilon$ and $\alpha - \epsilon$, making $J$ genuinely
non-smooth at isolated knife-edge points.  The FD check at L5 accordingly
shows relative error $\sim 10^{-2}$ (branch-switching noise) while the
same check at L4 with a gentler surface structure measures $10^{-6}$--$10^{-8}$
(the pipeline itself is exact).
\end{remark}

\subsection{Gauss-Newton recovery and the open research question}
\label{sec:e2recovery}

With the gradient $\partial J/\partial \alpha$ in hand, a Levenberg-Marquardt
Gauss-Newton step recovers the mode coefficients $\alpha^*$ from probe data.
The expected results at the sphere checkpoint ($k=4$ modes, L5):
\begin{itemize}
  \item Objective consistency: $J(\alpha^*) = O(10^{-30})$ (the sentinel ---
    the solution is represented exactly in the mode span).
  \item Recovery: $J$ drops $\sim 1000\times$ but \emph{plateaus} at a local
    minimum near $\alpha = 0$.  This is an honest, measured limitation: the
    objective has steep local structure at scales below $|\alpha^*|$
    (sensitivity at $0$ is $\sim 100\times$ the mean slope to the target),
    so descent methods stall.
\end{itemize}
The NS (Navier-Stokes) version of the same story (hero demo H1, steady
sphere) \emph{does} converge: advective transport enriches the measurement,
yielding recovery error $1.77 \times 10^{-3}$.  The M3 bunny headline
(\texttt{docs/dev/roadmap.md}) achieves recovery error $1.23 \times 10^{-3}$
via $h$/epoch continuation, confirming that Gauss-Newton alone loses to
landscape structure, while continuation across resolution and epoch beats it.

\subsection{Connection to E0 concepts}
\label{sec:e2e0connect}

\begin{itemize}
  \item \textbf{Do-not-differentiate list, item~3} (\S\ref{sec:dontdiff}):
    classification is frozen; the mesh arrays have zero derivative in $\alpha$.
  \item \textbf{IFT through the projection} (\S\ref{sec:dtd} and \S\ref{sec:e1ift}):
    the Newton projector is a relation, not an algorithm; the IFT closes the
    chain through the INR.
  \item \textbf{Identifiability / Gramian} (\S\ref{sec:identifiability}):
    rank-2 force aliasing (Gramian condition $\sim 80$) means two mode
    directions move the observable identically; the thick-band extraction and
    the Gramian check before any gradient run are the E2 equivalents of the
    beyond-FH gauge analysis.
  \item \textbf{Signal scale / probe placement} (\S\S\ref{sec:signalscale},
    \ref{sec:flatJ}): E2's 26-point shell at $r \approx 0.42$ is placed
    outside the wobble+edit envelope but well inside the Poisson decay radius,
    fixing the observability that corner probes (measured $J_0 \sim 4 \times
    10^{-4}$, wild Gauss-Newton) cannot provide.
  \item \textbf{The $2.5\times$ cost contract} (\S\ref{sec:cvm}): the
    adjoint for $k = 4$ mode coefficients costs one forward solve plus one
    transposed solve, regardless of $k$ --- the same $O(1)$ scaling as E0's
    material-parameter adjoint.
\end{itemize}

%-----------------------------------------------------------------------
\section{The Arc: E0 to E1/E2 and Beyond}
\label{sec:arc}

The E-track teaches a single idea in three registers of complexity.

\textbf{E0 (the on-ramp).} Build intuition for information-flow reversal
(\S\ref{sec:infoflow}), derive the discrete adjoint from the Lagrangian
(\S\ref{sec:lagrangian}), learn the taping rules (\S\S\ref{sec:tape},
\ref{sec:dontdiff}), handle transient chains (\S\ref{sec:transient}), and
design a good $J$ (Parts V--VI).  The payoff is the Allen-Cahn three-way
gradient check (\S\ref{sec:acpayoff}) and the beyond-FH gauge story
(\S\ref{sec:identifiability}).

\textbf{E1 (shape adjoint).} The same adjoint formula now acts on geometry:
the do-not-differentiate list items appear in the wild (classification
freeze), the custom VJP for the linear solver governs the adjoint solve,
and the IFT through the Newton projection closes the chain.  The gradient
cost is still one forward solve plus one transposed solve, regardless of
whether the shape is described by 3 parameters (disk) or 3000 (GridSDF
voxels).

\textbf{E2 (neural geometry).} GENIE's affine-in-$\alpha$ edit space turns
a neural SDF into a smooth, differentiable geometry oracle.  The framework
is unchanged: the tape collects SBM face cotangents, \texttt{torch.autograd}
propagates through the feature extractor, and the same adjoint solve closes
the loop.  The new content is the epoch-trust-region contract and the open
research question of landscape nonconvexity.

\textbf{Where the program goes.} The M3 bunny ladder
(\texttt{docs/dev/roadmap.md}) confirms that $h$/epoch continuation
converts an open research question (E2) into a headline result (recovery
error $1.23 \times 10^{-3}$): the machinery is complete; the open frontier
is richer observables and continuation strategies.  SP-1 (R3,
\texttt{docs/theory/\allowbreak structure\_property\_p3.md} and
\texttt{docs/dev/\allowbreak 2026-07-18-horizon-readiness.md}) builds the multi-protocol
inversion --- instrument-space $S(q,t)$, film height $h(t)$, PSF-convolved
images --- on the production stack, scaling to $60$--$140$~M dofs on Horizon
GB nodes.  The E-track is the on-ramp; the production inverse program is
the destination.

%=======================================================================
\begin{thebibliography}{9}

\bibitem{farrell2013}
P.~E.~Farrell, D.~A.~Ham, S.~W.~Funke, and M.~E.~Rognes,
\emph{Automated Derivation of the Adjoint of High-Level Mathematical
Programs},
SIAM Journal on Scientific Computing \textbf{35} (2013) C369--C393.
\url{https://doi.org/10.1137/120873558}

\bibitem{mitusch2019}
S.~K.~Mitusch, S.~W.~Funke, and J.~S.~Dokken,
\emph{dolfin-adjoint 2018.1: automated adjoints for FEniCS and Firedrake},
Journal of Open Source Software \textbf{4} (2019) 1292.
\url{https://doi.org/10.21105/joss.01292}

\bibitem{griewank2000}
A.~Griewank and A.~Walther,
\emph{Algorithm 799: Revolve: An Implementation of Checkpointing for the
Reverse or Adjoint Mode of Computational Differentiation},
ACM Transactions on Mathematical Software \textbf{26} (2000) 19--45.
\url{https://doi.org/10.1145/347837.347846}

\bibitem{beyondfh2026}
DiffSim group,
\emph{Beyond Flory-Huggins: differentiable $S(q)$ instrument-space
observable, gauge identifiability, and composition-coverage conditioning}
(internal dev note, 2026-07-14).
\texttt{docs/dev/2026-07-14-beyond-fh-learned-thermo.md}.

\bibitem{karki2025}
R.~Karki, V.~Krishnamurthy, and B.~Ganapathysubramanian,
\emph{GENIE: Editing Implicit Neural Representations via Last-Layer Edit
Space},
arXiv preprint arXiv:2603.29860 (2025).
\url{https://arxiv.org/abs/2603.29860}

\bibitem{neuralsdf2026}
DiffSim group,
\emph{NeuralSDF (Provided-INR) Oracle --- Implementation Spec (M1c/M2)}
(internal spec, 2026-07-06).
\texttt{docs/dev/specs/2026-07-06-neuralsdf.md}.

\end{thebibliography}

%=======================================================================
\appendix

\section{The Adjoint-Readiness Checklist}
\label{app:checklist}

\emph{Apply this checklist to any new PDE before writing gradient code.
The normative, citable version lives on the reference page
\texttt{docs/theory/}\linebreak\texttt{adjoint\_readiness\_checklist.md} (T4).
This appendix is the derived, annotated form; the reference page owns
the normative list.}

\medskip

\noindent\textbf{1. What to tape.}
Tape only the residual kernels that depend on the parameters you
differentiate with respect to.  Record $R(u, m)$ in $m$ (and in $u$ if
a later step or $J$ needs $\partial J/\partial u$).  Everything
algebraic-and-parameter-dependent is taped; the linear solver is not
(item~4).

\noindent\textbf{2. What to freeze.}
\begin{itemize}
\item Mesh, octree, element classification, preflight checks --- piecewise-constant
  in parameters (\S\ref{sec:dontdiff}, item~3).
\item Basis and quadrature tables: frozen inputs (\S\ref{sec:dontdiff}, item~4).
\item RNG seeds: reparametrize.
\item Any GP-field precomputed outside the tape: the M2 ``frozen at GPs''
  convention (\S\ref{sec:computeonce}, item~iii).
\end{itemize}

\noindent\textbf{3. What to store across steps.}
A transient adjoint sweeps backward and contracts with the primal state
$u^n$ at each step (\S\ref{sec:transient}).  Store the whole trajectory
(full-store, $O(N \times n)$ memory) or store checkpoints and recompute
segments (revolve, $O(\log N)$ memory at $O(N \log N)$ recompute cost).
Print the storage inventory before committing to a strategy.

\noindent\textbf{4. Transposed-solve infrastructure.}
The gradient path through any solver is the transposed solve at the
converged state (\S\S\ref{sec:lagrangian},\,\ref{sec:dontdiff}).  One LU
factorization serves both $A$ and $A^\top$ solves
(\S\ref{sec:computeonce}).  Build and test the adjoint solve before
writing the optimisation loop.

\noindent\textbf{5. The verification ladder.}
\begin{enumerate}
\item[(a)] Dot-product test $\langle Av,w\rangle = \langle v,A^\top w\rangle$
  per operator (\S\ref{sec:dotproduct}).
\item[(b)] Three-way gradient check: hand adjoint vs tape (or complex-step)
  vs central FD, for a small problem.  Agreement $< 10^{-6}$ relative
  error is the house standard.
\item[(c)] Transient-chain check: the full backward sweep vs FD end-to-end
  (\S\ref{sec:transient}).
\end{enumerate}
Climb the ladder before trusting a gradient in an optimiser.

\noindent\textbf{6. Nondifferentiability hazards.}
Replace $\mathrm{abs}$, $\min$, $\max$, threshold, and mask operations
with relaxed forms (\S\ref{sec:smoothness}): \texttt{softabs} ($\sqrt{x^2
+ \varepsilon^2}$), logsumexp \eqref{eq:lse}, sigmoid gates.  Do this
\emph{before} optimising.

\noindent\textbf{7. The $2.5\times$ cost contract.}
A correct reverse-mode adjoint adds $O(1)$ solves per forward step.
Measure the backward/forward wall-time ratio (\S\ref{sec:cvm}) before
scaling up.  If backward $\gg$ forward, something is being unrolled that
should not be.

\bigskip

\noindent\textbf{J-design addendum.}  Once the adjoint is verified:
\begin{itemize}
\item \emph{Signal scale first} (H4 protocol, \S\ref{sec:signalscale}):
  confirm $|\partial \hat{J}/\partial m|$ clears the FD noise floor by
  $\geq 10^3$ before any optimisation.
\item \emph{Probe placement matters} (\S\ref{sec:flatJ}): an observable
  blind to the active region gives a vanishing gradient; design the
  measurement to see the physics.
\item \emph{Identifiability first} (\S\ref{sec:identifiability}): compute
  the observation-Jacobian condition number before fitting.  Remove gauge
  freedoms from the parameterisation by construction.  Add Tikhonov
  regularisation as a scientific prior, not as a solver fix.
\item \emph{Diverse observables} (\S\ref{sec:multiobs}): combine
  probe misfits, structure-factor misfits, and integral observables to
  illuminate different directions of parameter space.
\end{itemize}

\end{document}
